% Modernised reading — fonds Alexandre Grothendieck, shelfmark 94
% (pages 1-23, the whole folder).
%
% This is an INTERPRETATION, not a transcription. It re-reads the pages in
% today's notation and vocabulary, states the hypotheses the manuscript leaves
% implicit, and drops the critical apparatus so that the mathematics reads
% straight through. Everything the brackets and underlines carried moves into
% a footnote. It is held to being mathematically correct as it stands; where
% the manuscript is loose or mistaken, the true statement is what appears here
% and the footnote says what the page has. In any dispute of substance the
% transcription governs: whoever cites Grothendieck cites batch-01.fr.tex
% (pp. 1-20) or batch-02.fr.tex (pp. 21-23).
%
% Pass: Opus 5.5 (claude-opus-5-5), 2026-10-10 — first pass, unchecked against the pages by a human.
% The folder was read whole, its two batches together; this is one document
% for the shelfmark. It was made on Opus 5.5 and has not been measured against
% the reference reading of folder 115 (made on Opus 5). The literature named
% in the footnotes is cited from memory and was not checked against the
% sources.
%
% Spine. One numbered run (statements 1 to 21), under the cover title
% « Construction des Modules cohérents par leurs foncteurs représentatifs »:
%   I    pp. 2-4    descent and completion: Lemme 1, Lemme 2, Prop. 3.
%   II   pp. 4-7    the local case and pro-representability: Prop. 4,
%                   Lemme 5, Remarque 5 bis.
%   III  pp. 8-15   a test module: Tor-vanishing against regular sequences
%                   (Lemme 6, Cor. 7), the « faisceau fondamental » (Lemme 8,
%                   Cor. 9, 10), stationarity of a pro-system (Cor. 11).
%   IV   pp. 15-16  Théorème 12, the criterion; Remarque 13 (interrupted).
%   V    pp. 17-20  Théorème 14: Hom(P, M) = f_* Hom(F, G (x) M), with the
%                   generic base-change Lemmes 15-17.
%   VI   pp. 21-23  Remarque 18 (« foncteurs constructibles »), Théorème 19
%                   (Weil restriction of affine X-schemes along a proper flat
%                   f.p. X -> S), Cor. 20, Remarque 21.
% Pages 11 and 13 are typed leaves unrelated to the notes (skipped by the
% transcription); page 1 is the cover.
%
% Notation fixed for the folder. A noetherian ring (except in Lemme 1);
% F : Mod(A) -> (Ens) a covariant functor, « représentable » meaning
% isomorphic to Hom_A(P, -); F_B for its restriction to B-modules; F_g, F_p
% for B = A_g, A_p. The page writes f for the element of A, the morphism and
% the regular sequences; here g in A, f : X -> S, t = (t_1, ..., t_n) for a
% regular sequence or a system of parameters. Completions: \widehat{A}^{(g)}
% the g-adic one (pp. 2-3, 15), \widehat{A}_p the completion of the local
% ring A_p (pp. 4-5, 15) — the page writes both \widehat{A}. Sheaves on X in
% sections V-VI are calligraphic (\mathcal F, \mathcal G), so as not to clash
% with the functor F. Ω is a canonical module (his « faisceau fondamental »).
%
% Substantive departures from the pages, with pages.
%  - p. 2, Lemme 2 (ii): stated on the page for completions of A-modules of
%    finite type; the proof (p. 3) needs it for every finite Â-module, which
%    is the form Prop. 3 (iii) gives. Stated in that form.
%  - p. 4, Prop. 4 (iv): the parenthesis is partly illegible; read « it
%    suffices for M of finite length », which is what the proof uses.
%  - p. 6: the gluing step of Lemme 5 is largely illegible; the reading gives
%    a complete argument (Nakayama along V(p)) built on the legible fragments,
%    footnoted as ours.
%  - p. 7, Rem. 5 bis: the converse is only sketched and partly illegible;
%    its structure is given, not a proof.
%  - p. 9, Cor. 7: the letter read Q is Ω; « Z^1(K) » is K^1 for a two-term
%    complex; the argument goes through with all higher Tor (ours).
%  - p. 10, Cor. 10: Ω put on U, not X; no proof on the page.
%  - p. 14, Cor. 11: the system (P_i) must be strict, of modules of finite
%    type (supplied; the cancelled first state has « strict »).
%  - p. 16, Rem. 13: interrupted. (vi) => (v) is justified here; the claim
%    that the « quotient of a regular ring » hypothesis becomes useless is
%    clear only for integral quotients and is left open. Rem. 5 bis's
%    hypothesis on normal loci is not checked by Th. 12 for a mere quotient of
%    a regular ring; it holds in the application (A of finite type over Z).
%  - p. 19, Lemme 16: gr_J(O_X) read gr_J(R); the formula's F read R.
%  - p. 22, 1st case: « th. 12 » is Th. 14 with G = O_X.
%
% Modern names given and footnoted as ours: (co)representable functor,
% faithfully flat descent, Matlis duality, canonical module, generic flatness,
% projection formula, coherent functors (Auslander, Hartshorne), Weil
% restriction; Th. 14 is in EGA III_2 § 7.7 for a noetherian base.
%
% What the folder announces and never establishes: the proof of Cor. 10;
% the converse of Rem. 5 bis in detail; the end of Rem. 13; the « NB » of
% p. 17 (W(...) for G quasi-coherent); the programme of Rem. 18. « l'exposé
% Murre, th. 3 » (p. 23) is cited, not proved.
%
% Group. The folder sits in the inventory group « Géométrie et topologie
% combinatoire » (69-102); its subject does not match (as for 90, 100, 101).
% It meets none of the combinatorial readings taken as models (81, 88, 89).
% It meets 101 (the « module fondamental », pp. 18-20 of 101) and 100 (Weil
% restriction, pp. 47-50 of 100); the heading « Cor. Main » is read as in the
% transcription of 90 (pp. 5, 9).
\documentclass[11pt,a4paper]{article}
\input{../preamble/grothendieck}

\folder{94}
\pages{1}{23}
\dating{[à partir de 1965]}
\watermark{Édition de démonstration\\interprétation personnelle de l'œuvre}
\foldertitle{Dualité des modules / Construction de faisceaux de modules : notes manuscrites (s.d.) — lecture modernisée du dossier entier}

\begin{document}

\section*{Résumé}

\begin{resume}
Ces feuillets cherchent à construire un module --- un objet de l'algèbre
commutative, qui généralise l'espace vectoriel --- sans jamais le voir
directement, en ne connaissant que la manière dont il s'envoie dans tous les
autres. Une clé se reconnaît aux serrures qu'elle ouvre : si l'on sait, pour
chaque serrure, si elle s'ouvre et comment, on connaît la clé. De même, un
module $P$ est entièrement déterminé par la règle qui associe à chaque module
$M$ l'ensemble $\mathrm{Hom}(P, M)$ des applications linéaires de $P$ dans $M$.
C'est ce que la couverture appelle le « foncteur représentatif » de $P$.

La question du dossier est la question inverse. On se donne une règle
$M \mapsto F(M)$, qui à chaque module associe un ensemble, de façon compatible
aux applications linéaires. Quand cette règle est-elle celle d'une clé ?
Autrement dit, quand existe-t-il un module $P$, de type fini, tel que
$F(M) = \mathrm{Hom}(P, M)$ pour tout $M$ ? La question se pose parce qu'en
géométrie algébrique les objets qu'on veut construire --- un espace qui
paramètre des morphismes, des sections, des déformations --- se présentent
d'abord ainsi : on sait dire ce qu'ils « devraient » être vus de chaque test,
et il faut prouver qu'il y a quelque chose derrière.

La réponse est une liste de conditions, et chacune a un sens intuitif : $F$
doit respecter les recollements (il est exact à gauche), les réunions
croissantes (il commute aux limites inductives filtrantes), les passages à la
limite par approximations successives (une condition de complétion), il doit
prendre des valeurs finies sur des arguments finis, et il doit se comporter
génériquement de façon régulière. Le dernier ingrédient est le plus fin : pour
savoir si une suite de modules de plus en plus précis finit par se stabiliser,
il suffit de la tester contre un seul module bien choisi, le « faisceau
fondamental », que l'on appelle aujourd'hui module canonique. C'est le seul
endroit où la dualité du titre du dossier intervient : comme on reconnaît un
espace vectoriel de dimension finie à son dual, on reconnaît ici un module à
ses applications vers ce module témoin.

Le critère (théorème 12) est ensuite appliqué deux fois. D'abord à une famille
de variétés projectives paramétrée par une base $S$ : les morphismes entre
deux faisceaux sur les fibres, qui varient avec le point de la base, sont
« représentés » par un seul module $P$ sur $S$ (théorème 14). Puis à la
construction d'espaces de sections : si $Z$ est un objet affine au-dessus de
$X$, l'ensemble des sections de $Z$ sur les fibres de $X \to S$ est lui-même un
schéma affine sur $S$ (théorème 19). En chemin, les feuillets proposent de
baptiser « constructibles » les foncteurs qui satisfont à presque toutes ces
conditions, et esquissent le programme d'une théorie.

Pour aller plus loin : foncteurs représentables et pro-représentables, descente
fidèlement plate, dualité de Matlis et module canonique, platitude générique et
changement de base en cohomologie, foncteurs cohérents au sens d'Auslander,
restriction de Weil.

\keywords{representable functor, pro-representable functor, faithfully flat descent, Matlis duality, canonical module, generic flatness, cohomology and base change, coherent functor, Weil restriction}
\end{resume}

\section*{Le fil du dossier, et les conventions}

\subsection*{Les stations}

\begin{enumerate}
\item[I.] \emph{Descente et complétion} (pages 2 à 4) : on ramène la
représentabilité à celle sur un ouvert $\operatorname{Spec} A_g$ et sur le
complété $g$-adique (lemmes 1 et 2), puis, par récurrence noethérienne, à une
condition générique (proposition 3).
\item[II.] \emph{Le cas local, la pro-représentabilité} (pages 4 à 7) : sur un
anneau local, une condition de finitude testée sur l'enveloppe injective
suffit (proposition 4) ; globalement, $F$ est pro-représentable par un système
qui se stabilise en chaque point (lemme 5), et la condition générique peut se
remplacer par la constructibilité d'une fonction rang (remarque 5 bis).
\item[III.] \emph{Un module témoin} (pages 8 à 15) : annulation générique des
$\mathrm{Tor}$ contre les suites régulières (lemme 6, corollaire 7), le
faisceau fondamental et ses réductions (lemme 8, corollaires 9 et 10), et la
stabilisation d'un pro-système testée sur lui (corollaire 11).
\item[IV.] \emph{Le critère} (pages 15 et 16) : théorème 12, et la
remarque 13, interrompue.
\item[V.] \emph{Le module $P$ d'un morphisme propre} (pages 17 à 20) :
théorème 14, et les lemmes 15 à 17 de changement de base générique.
\item[VI.] \emph{Foncteurs constructibles, restriction de Weil} (pages 21 à
23) : remarque 18, théorème 19, corollaire 20, remarque 21.
\end{enumerate}

Les énoncés sont numérotés de 1 à 21 d'un bout à l'autre ; le dossier est
d'une seule coulée. Le titre de la couverture --- « Construction des Modules
cohérents par leurs foncteurs représentatifs. Application à la
représentabilité d'un foncteur [\ldots] d'une section, à la définition des
invariants $\Omega^{1}(y)$\,\ldots » --- en est le programme : la première
phrase est le théorème 12, la seconde le théorème 19\footnote{Le titre est
écrit à l'encre dans l'angle d'une chemise brune (page 1), avec « et de schémas
finis sur d'autres par descente » biffé ; deux mots après « foncteur $G/S$ »
sont illisibles, et « $\Omega^{1}(y)$ » est une lecture incertaine. Rien dans
le dossier ne revient à ces invariants. Aucune page n'est intitulée « Dualité
des modules », la première moitié du titre de l'inventaire ; la dualité n'y
paraît que comme outil, aux pages 8 à 15. Les pages 11 et 13 sont des
feuillets dactylographiés sans rapport, qui ne portent pas sa main.}.

\subsection*{Conventions, valables pour tout le dossier}

$A$ est un anneau commutatif, noethérien sauf au lemme 1. Un \emph{foncteur}
est un foncteur covariant $F \colon \mathrm{Mod}(A) \to (\mathrm{Ens})$. Il est
\emph{représentable} s'il est isomorphe à $\mathrm{Hom}_A(P, -)$ pour un
$A$-module $P$ --- on dit aussi, pour un foncteur covariant, coreprésentable ;
c'est l'usage de ces pages qu'on garde --- et $P$ est alors unique à
isomorphisme unique près. Il est \emph{pro-représentable} s'il est isomorphe à
$M \mapsto \varinjlim_i \mathrm{Hom}(P_i, M)$ pour un système projectif
filtrant $(P_i)$, \emph{strict} si les morphismes de transition sont
surjectifs. \emph{Exact à gauche} veut dire : commute aux limites projectives
finies (noyaux, produits finis).

Pour une $A$-algèbre $B$, $F_B$ est le composé de la restriction des scalaires
$\mathrm{Mod}(B) \to \mathrm{Mod}(A)$ et de $F$ ; si $F = \mathrm{Hom}_A(P, -)$,
alors $F_B = \mathrm{Hom}_B(P \otimes_A B, -)$. On écrit $F_g$ et
$F_{\mathfrak{p}}$ pour $B = A_g$ et $B = A_{\mathfrak{p}}$.

Deux complétions interviennent, que la page écrit toutes deux
$\widehat{A}$ : $\widehat{A}^{(g)} = \varprojlim A/g^{n+1}A$, complété
$g$-adique, et $\widehat{A}_{\mathfrak{p}}$, complété de l'anneau local
$A_{\mathfrak{p}}$, d'idéal maximal $\mathfrak{m}_{\mathfrak{p}}$. Les éléments
de $A$ sont notés $g$, les suites régulières et systèmes de paramètres
$\underline{t} = (t_1, \ldots, t_n)$, et $f$ est réservé aux morphismes de
schémas\footnote{La page écrit $f$ pour les trois. On a changé deux d'entre
eux pour que les sections V et VI, où les trois se rencontrent, restent
lisibles.}. Aux sections V et VI, les faisceaux sur $X$ sont notés
$\mathcal{F}$, $\mathcal{G}$, la lettre $F$ restant au foncteur.

Une observation, que les pages utilisent sans la dire, sert partout : \emph{si
$F$ est exact à gauche, chaque $F(M)$ est un $A$-module} (l'addition vient de
$F(M \times M) = F(M) \times F(M)$, les scalaires de $F$ appliqué aux
homothéties), \emph{et si $F$ commute de plus aux limites inductives
filtrantes, $F(M_g) = F(M)_g$ et $F(M_{\mathfrak{p}}) = F(M)_{\mathfrak{p}}$},
puisque $M_g$ est la limite de $M \xrightarrow{g} M \xrightarrow{g} \cdots$.
En particulier un élément de $F(M)$ est nul s'il l'est en chaque point.

\subsection*{Ce que le dossier annonce sans l'établir}

La démonstration du corollaire 10 (existence d'un module canonique sur un
ouvert dense), le détail de la réciproque de la remarque 5 bis, la fin de la
remarque 13, la note marginale de la page 17, le programme de la remarque 18.
Le théorème de « l'exposé Murre » (page 23) est invoqué, non démontré.

\pagerange{2}{4}
\section*{I. Descente et complétion (pages 2 à 4)}

\textbf{Lemme 1.} Soient $A$ un anneau quelconque et $B$ une $A$-algèbre
fidèlement plate. Pour que $F$ soit représentable (resp.\ par un module de type
fini, de présentation finie), il faut et il suffit que $F$ soit exact à gauche
et que $F_B$ soit représentable (resp.\ par un $B$-module de type fini, de
présentation finie).

La page ne donne pas la démonstration ; c'est la descente fidèlement plate.
Si $F_B$ est représenté par $Q$, les deux restrictions de $F_{B \otimes_A B}$
sont représentées par $Q \otimes_A B$ et $B \otimes_A Q$, d'où par Yoneda une
donnée de descente sur $Q$, qui descend en un $A$-module $P$ ; l'exactitude à
gauche de $F$ appliquée à la suite exacte
$M \to M \otimes_A B \rightrightarrows M \otimes_A B \otimes_A B$ montre que $P$
représente $F$, et les propriétés de finitude descendent\footnote{La
démonstration est de cette lecture. Sur la page, « Supposons $F$ exact à
gauche, et soit » est biffé et remplacé par « Soit $B$ une $A$-algèbre fid.\
plate » ; l'exactitude à gauche passe dans la conclusion.}.

\textbf{Lemme 2.} Soient $A$ noethérien et $g \in A$. Pour que $F$ soit
représentable par un module de type fini, il faut et il suffit que :
\begin{enumerate}
\item[(i)] $F$ soit exact à gauche ;
\item[(ii)] pour tout $\widehat{A}^{(g)}$-module de type fini $M$,
$F(M) \xrightarrow{\ \sim\ } \varprojlim_n F(M/g^{n+1}M)$ ;
\item[(iii)] $F$ commute aux limites inductives filtrantes ;
\item[(iv)] $F_g$ soit représentable ;
\item[(v)] pour tout $n$, $F_{A/g^{n}A}$ soit représentable.
\end{enumerate}
La nécessité est claire ; (iii) traduit le fait qu'un module $P$ tel que
$\mathrm{Hom}(P, -)$ commute aux limites inductives filtrantes est de
présentation finie, donc de type fini\footnote{La page énonce (ii) pour les
complétés $\widehat{M}$ des $A$-modules de type fini $M$. Sous cette forme la
condition est plus faible que ce que la démonstration utilise : il faut
connaître $F$ sur le module $\widehat{P}$ construit plus bas, qui n'a pas de
raison d'être le complété d'un $A$-module. La forme retenue ici est celle de la
proposition 3 (iii), où la page corrige d'ailleurs $F(\widehat{M})$ en
$F(M)$ ; elle est encore nécessaire, puisqu'un tel $M$ est séparé et complet.
« Lemme 2 » surcharge un « Corollaire 2 » biffé.}.

\emph{Suffisance.} L'algèbre $B = \widehat{A}^{(g)} \times A_g$ est fidèlement
plate sur $A$ : elle est plate, et son spectre recouvre
$V(g) \cup D(g)$. Comme $F$ est exact à gauche,
$F_B(N_1 \times N_2) = F(N_1) \times F(N_2)$, et il suffit par le lemme 1 que
$F_g$ --- c'est (iv) --- et $F_{\widehat{A}^{(g)}}$ soient représentables ; le
$P$ obtenu sera de type fini par (iii). Posons $A_n = A/g^{n+1}A$. Par (v),
$F_{A_n}$ est représenté par un $A_n$-module $P_n$, de type fini par (iii) ;
par unicité, $P_{n} \otimes A_{n-1} = P_{n-1}$, les $P_n$ « se recollent », et
$\widehat{P} = \varprojlim P_n$ est un $\widehat{A}^{(g)}$-module de type fini.
Pour $M$ de type fini sur $\widehat{A}^{(g)}$, (ii) donne
\[
  F(M) = \varprojlim F(M/g^{n+1}M) = \varprojlim
  \mathrm{Hom}(P_n, M/g^{n+1}M) = \mathrm{Hom}(\widehat{P}, M),
\]
et, tout $\widehat{A}^{(g)}$-module étant réunion filtrante de sous-modules de
type fini, (iii) étend l'isomorphisme à tous les modules.

\textbf{Proposition 3.} Soit $A$ noethérien. Pour que $F$ soit représentable
par un module de type fini, il faut et il suffit que :
\begin{enumerate}
\item[(i)] $F$ soit exact à gauche ;
\item[(ii)] $F$ commute aux limites inductives filtrantes\footnote{La page
précise, en interligne et d'une lecture incertaine, « monomorphes » : limites
inductives filtrantes à morphismes de transition injectifs. C'est en effet
tout ce que la démonstration du lemme 2 utilise sur un anneau noethérien
--- tout module est la réunion filtrante de ses sous-modules de type fini, et
$P$ est de type fini dès que $\mathrm{Hom}(P,-)$ commute à ces réunions.} ;
\item[(iii)] pour tout $g \in A$ et tout $\widehat{A}^{(g)}$-module de type fini
$M$, $F(M) \xrightarrow{\ \sim\ } \varprojlim F(M/g^{n+1}M)$ ;
\item[(iv)] pour tout quotient $B \neq 0$ de $A$, il existe $g \in B$ non
nilpotent tel que $F_{B_g}$ soit représentable.
\end{enumerate}

\emph{Démonstration.} Les conditions passent de $F$ à $F_B$ pour tout quotient
$B$ de $A$. Par récurrence noethérienne, on peut supposer $F_{A/J}$
représentable pour tout idéal $J \neq 0$, et $A \neq 0$. Soit $g$ comme en
(iv) pour $B = A$ ; $g$ n'étant pas nilpotent, $g^{n} \neq 0$ et $F_{A/g^{n}A}$
est représentable pour tout $n$. Le lemme 2 conclut\footnote{Sur la page, la
condition (iv) est écrite avant (iii), d'abord numérotée (iii), et une flèche
la renvoie à sa place.}.

« En pratique, la condition (iv) est parfois difficile à vérifier » : tout ce
qui suit, jusqu'au théorème 12, sert à la remplacer.

\pagerange{4}{7}
\section*{II. Le cas local et la pro-représentabilité (pages 4 à 7)}

\pagerange{4}{5}
\subsection*{L'anneau local (pages 4 et 5)}

\textbf{Proposition 4.} Soit $A$ local noethérien, d'idéal maximal
$\mathfrak{m}$, de complété $\widehat{A}$. Pour que $F$ soit représentable par
un module de type fini, il faut et il suffit que :
\begin{enumerate}
\item[(i)] $F$ soit exact à gauche ;
\item[(ii)] $F$ commute aux limites inductives filtrantes ;
\item[(iii)] pour tout $\widehat{A}$-module de type fini $M$,
$F(M) \xrightarrow{\ \sim\ } \varprojlim F(M/\mathfrak{m}^{n+1}M)$ ;
\item[(iv)] pour tout $\widehat{A}$-module de type fini $M$, $F(M)$ soit de
type fini sur $\widehat{A}$ ; il suffit de le savoir pour $M$ de longueur
finie\footnote{La parenthèse de (iv) est en partie illisible ; on la lit « il
suffit [de le savoir] si $M$ est [de longueur] fini[e] », qui est ce que la
démonstration utilise. La page écrit $F(\widehat{M})$ en (iii) et (iv), le
chapeau ajouté après coup ; pour un $\widehat{A}$-module de type fini,
$\widehat{M} = M$.}.
\end{enumerate}

\emph{Démonstration.} Comme $\widehat{A}$ est fidèlement plat sur $A$, le
lemme 1 permet de supposer $A$ complet. Posons $A_n = A/\mathfrak{m}^{n+1}$.
Sur les $A_n$-modules de longueur finie, catégorie artinienne, le foncteur
exact à gauche $F$ est pro-représentable, par un système strict
$(Q_i)$\footnote{C'est le théorème de pro-représentabilité de Grothendieck
(« Technique de descente et théorèmes d'existence en géométrie algébrique,
II », Séminaire Bourbaki, exposé 195, 1960) ; la page dit seulement « est
pro-représentable ».}. Soit $I_n$ l'enveloppe injective du corps résiduel sur
$A_n$ --- « l'injectif indécomposable de $A_n$ » ---, de longueur finie. Alors
$F(I_n) = \varinjlim \mathrm{Hom}(Q_i, I_n)$ est une réunion croissante des
duaux de Matlis des $Q_i$ ; elle est de longueur finie par (iv), donc la suite
se stabilise, et, la dualité de Matlis étant fidèle, le système $(Q_i)$ aussi :
$F_{A_n}$ est représenté par un $A_n$-module $P_n$ de type fini\footnote{Le
raisonnement par dualité de Matlis est de cette lecture ; la page dit « en
vertu de la condition (iv), testée avec l'injectif indécomposable $I_n$ de
$A_n$, on trouve qu'il est représentable ».}. Les $P_n$ se recollent en
$P = \varprojlim P_n$, de type fini ; (iii) montre que $P$ représente $F$ sur
les modules de type fini, et (ii) sur tous.

\pagerange{5}{6}
\subsection*{Pro-représentabilité (pages 5 et 6)}

\textbf{Lemme 5.} Soit $A$ noethérien, $S = \operatorname{Spec} A$, et
supposons :
\begin{enumerate}
\item[(i)] $F$ exact à gauche ;
\item[(ii)] $F$ commute aux limites inductives filtrantes ;
\item[(iii)] pour tout $\mathfrak{p} \in S$ et tout
$\widehat{A}_{\mathfrak{p}}$-module de type fini $M$,
$F(M) \xrightarrow{\ \sim\ } \varprojlim F(M/\mathfrak{m}_{\mathfrak{p}}^{n+1}M)$ ;
\item[(iv)] pour tout $\mathfrak{p}$ et tout $A_{\mathfrak{p}}$-module de
longueur finie $M$, $F(M)$ est un $A_{\mathfrak{p}}$-module de type fini ;
\item[(v)] pour tout idéal premier $\mathfrak{p}$, il existe
$g \in A/\mathfrak{p}$ non nul tel que $F_{(A/\mathfrak{p})_g}$ soit
représentable.
\end{enumerate}
Alors $F$ est pro-représentable par un système projectif strict $(P_i)$ de
modules de type fini, et pour tout $\mathfrak{p}$ le système des
$(P_i)_{\mathfrak{p}}$ est stationnaire : $F_{\mathfrak{p}}$ est représentable
par un $A_{\mathfrak{p}}$-module de type fini\footnote{Sur la page, la
conclusion est écrite d'abord, et un long trait la renvoie après les
conditions ; « Lemme 5 » surcharge « Proposition », biffé. En (v), les
$A/\mathfrak{p}$ sont écrits sur des $A_{\mathfrak{p}}$.}.

\emph{Démonstration.} Pour chaque $\mathfrak{p}$, la proposition 4 appliquée à
$A_{\mathfrak{p}}$ montre, par (i) à (iv), que $F_{\mathfrak{p}}$ est
représenté par un $A_{\mathfrak{p}}$-module de type fini $P_{\mathfrak{p}}$ ;
par unicité, $(P_{\mathfrak{p}})_{\mathfrak{q}} = P_{\mathfrak{q}}$ pour
$\mathfrak{q} \subset \mathfrak{p}$ --- les $P_{\mathfrak{p}}$ sont
« compatibles par générisation ».

Compte tenu de (i) et (ii), il suffit, pour la pro-représentabilité, que tout
$\alpha \in F(M)$, $M$ de type fini, ait un \emph{module de définition} : un
plus petit sous-module $N \subset M$ tel que $\alpha \in F(N)$. En chaque point
on en a un : $N_{\mathfrak{p}} = \operatorname{Im}(\alpha_{\mathfrak{p}} \colon
P_{\mathfrak{p}} \to M_{\mathfrak{p}})$. Il suffit que la famille des
$N_{\mathfrak{p}}$ provienne d'un sous-module $N$ de $M$ ; celui-ci convient
alors, l'appartenance à $F(N) = \ker\bigl(F(M) \to F(M/N)\bigr)$ se vérifiant
point par point.

Soit $U \subset S$ le plus grand ouvert sur lequel les $N_{\mathfrak{p}}$ se
recollent en un sous-faisceau cohérent $\underline{N}_U \subset
\widetilde{M}|_U$, et supposons $U \neq S$. Soit $x = \mathfrak{p}$ un point
maximal de $S - U$. Les générisations strictes de $x$ sont dans $U$, et
$\widetilde{N_{\mathfrak{p}}}$ coïncide avec $\underline{N}_U$ sur
$U \cap \operatorname{Spec} A_{\mathfrak{p}}$ ; on peut donc les recoller sur
un voisinage ouvert $V$ de $x$ : il existe $\underline{N}' \subset
\widetilde{M}|_V$ cohérent avec $\underline{N}'_x = N_x$ et
$\underline{N}'_s = N_s$ pour $s \in U \cap V$. Reste à voir que, $V$ étant
assez petit, $\underline{N}'_s = N_s$ aussi pour $s \in V - U$, qu'on peut
supposer égal à $V \cap V(\mathfrak{p})$.

D'abord, $\alpha_x$ appartenant à $F(N_x) = \varinjlim F(\underline{N}'(D(h)))$,
on peut, quitte à restreindre $V$, supposer $\alpha|_V \in F(\underline{N}')$ ;
alors $N_s \subset \underline{N}'_s$ pour tout $s \in V$, et l'on est ramené au
cas $\underline{N}' = \widetilde{M}|_V$. Ensuite, par (v), on peut prendre
$V = \operatorname{Spec} A_h$ avec $F_{(A/\mathfrak{p})_h}$ représenté par un
module de type fini $Q$. L'image de $\bar{\alpha} \colon Q \to
(M/\mathfrak{p}M)_h$ est tout au point générique $x$ --- car $N_x = M_x$ ---,
donc, $V$ restreint encore, partout. Pour $s \in V \cap V(\mathfrak{p})$,
$Q_s = P_s \otimes A_s/\mathfrak{p}A_s$ (les deux représentent le même
foncteur), donc $N_s + \mathfrak{p}M_s = M_s$, et comme $\mathfrak{p} \subset
\mathfrak{m}_s$, Nakayama donne $N_s = M_s$. Ainsi les $N_{\mathfrak{p}}$ se
recollent sur $U \cup V \supsetneq U$, contradiction\footnote{Ce paragraphe
et le précédent reconstruisent la page 6, dont la fin est très remaniée et
pour une bonne part illisible : l'ordre des fragments est incertain. Ce qui se
lit --- le point maximal $x$ de $S - U$ avec « N.B.\ $x = \mathfrak{p}$ ! »,
le recollement sur un voisinage $V$, le choix de $V = \operatorname{Spec} A_h$
« grâce à la condition (v) », « $\alpha_V \in F(\underline{N}')$ donc $N_s
\subset \underline{N}'_s$ », une image « tout entière » le long de
$V(\mathfrak{p})$ « par Nakayama », et en marge « Donc on est ramené au cas
$\underline{N}' = M$ » --- est ce que la démonstration ci-dessus enchaîne ;
l'usage précis du module $Q$ est de cette lecture.}.

\pagerange{7}{7}
\subsection*{Une condition de constructibilité (page 7)}

\textbf{Remarque 5 bis.} Supposons que, pour tout quotient intègre $B$ de $A$,
le lieu normal de $\operatorname{Spec} B$ contienne un ouvert non
vide\footnote{La page écrit « $\mathrm{Nor}(B)$ contient un ouvert non vide »,
lecture incertaine, ajoutée en interligne : il faut l'entendre du lieu normal.
Cette condition est satisfaite si $A$ est excellent --- par exemple de type
fini sur $\mathbb{Z}$ ou sur un corps ---, puisque le lieu normal y est alors
ouvert et contient le point générique.}. Alors, en présence de (i) à (iv), la
condition (v) du lemme 5 équivaut à :
\begin{enumerate}
\item[(v bis)] la fonction $\mathfrak{p} \mapsto
\dim_{k(\mathfrak{p})} F(k(\mathfrak{p}))$ est constructible sur
$\operatorname{Spec} A$.
\end{enumerate}

Que (v) entraîne (v bis) est « trivial » : sur $\operatorname{Spec}
(A/\mathfrak{p})_g$, $F(k(\mathfrak{q})) = \mathrm{Hom}(Q, k(\mathfrak{q}))$ a
pour dimension le rang de la fibre de $Q$, constant sur un ouvert non vide de
$V(\mathfrak{p})$ ; et une fonction sur un espace noethérien qui est constante
sur un ouvert non vide de chaque fermé irréductible est constructible.

La réciproque n'est qu'esquissée, et le milieu de la page se lit mal. Sa
structure est la suivante. Pour $A$ intègre, la fonction rang est constante,
égale à $n$, sur un ouvert non vide $U$ ; sur un anneau réduit, un module de
type fini dont les fibres ont toutes même rang est libre, de sorte que les
$P_{\mathfrak{p}}$, $\mathfrak{p} \in U$, sont libres de rang $n$, et forment
des réseaux d'un même espace vectoriel sur le corps des fractions. On les
place, grâce à la commutation aux limites inductives, dans un module de type
fini $N$, et la finitude de $F(N)$ donne $P_{\mathfrak{p}} = N_{\mathfrak{p}}$
en tous les points de codimension 1 de $U$, sauf un nombre fini, qu'on
retire. Si enfin $U$ est normal, deux réseaux libres qui coïncident en
codimension 1 coïncident partout, et $N$ représente $F$ sur
$U$\footnote{Comment $N$ est obtenu, et pourquoi la coïncidence en
codimension 1 vaut hors d'un ensemble fini, se trouve dans les lignes
illisibles ; nous ne le reconstituons pas. L'argument final --- un module
libre sur un anneau normal est l'intersection de ses localisés en
codimension 1 --- est l'extension de Hartogs, et c'est lui qui demande
l'hypothèse sur le lieu normal.}.

\pagerange{8}{15}
\section*{III. Un module témoin (pages 8 à 15)}

Il s'agit maintenant de vérifier la condition (v) du lemme 5 sans la
supposer : trouver, génériquement, un module contre lequel tester la
stabilisation du système $(P_i)$. Les outils sont l'annulation générique des
$\mathrm{Tor}$ contre les suites régulières, puis le module canonique.

\pagerange{8}{9}
\subsection*{Annulation générique des $\mathrm{Tor}$ (pages 8 et 9)}

\textbf{Lemme 6.} Soient $S$ un schéma localement noethérien, $\mathcal{J}$ son
nilradical, $S_0 = S_{\mathrm{red}}$, et $R$ un $\mathcal{O}_S$-Module cohérent.
\begin{enumerate}
\item[a)] Il existe un ouvert dense $U$ de $S$ sur lequel
$\mathrm{gr}_{\mathcal{J}}(\mathcal{O}_S)$ et $\mathrm{gr}_{\mathcal{J}}(R)$
sont des $\mathcal{O}_{S_0}$-Modules localement libres.
\item[b)] Pour un tel $U$, tout $s \in U$ et toute suite
$\mathcal{O}_{S,s}$-régulière $\underline{t}$ d'éléments de $\mathfrak{m}_s$,
\[
  \mathrm{Tor}_i^{\mathcal{O}_{S,s}}\bigl(R_s,\ \mathcal{O}_{S,s}/(\underline{t})\bigr) = 0
  \qquad (i > 0) .
\]
\end{enumerate}

\emph{Démonstration.} a) est la liberté générique d'un Module cohérent sur un
schéma réduit. Pour b), posons $A = \mathcal{O}_{S,s}$, $J$ son nilradical,
$A_0 = A/J$. Le complexe de Koszul $K_{\cdot}(\underline{t})$ résout
$A/(\underline{t})$, et il suffit que $\underline{t}$ soit $R_s$-régulière.
On montre, par récurrence sur $n$, que les images $\bar{t}_i$ forment une suite
régulière de $A_0$ et que $\mathrm{gr}_J(R_s/(\underline{t})R_s) =
\mathrm{gr}_J(R_s)/(\bar{\underline{t}})$ est libre sur
$A_0/(\bar{\underline{t}})$ ; un élément régulier sur chaque gradué d'une
filtration finie est régulier sur le module. Pour l'étape, si $\bar{t}_{k+1}$
était diviseur de zéro dans $A_0/(\bar{t}_1, \ldots, \bar{t}_k)$, il le serait
dans le dernier gradué non nul de $A/(t_1, \ldots, t_k)$, qui en est un
sous-module libre non nul, contre la régularité de $t_{k+1}$\footnote{La page
dit « c'est bien standard si $n = 1$, et on continue par récurrence » ; le
pas de récurrence est de cette lecture. Le $M$ de « suite $M$-régulière » n'y
est défini que dans une ligne biffée de l'énoncé ; c'est le module $R_s$ (que
la page appelle $F$, nom qu'on réserve au foncteur).}.

\textbf{Corollaire 7.} Soient $S$ noethérien, $P$ et $\Omega$ des
$\mathcal{O}_S$-Modules cohérents. Il existe un ouvert dense $U$ de $S$ tel que,
pour tout $s \in U$ et toute suite $\mathcal{O}_{S,s}$-régulière
$\underline{t}$ de $\mathfrak{m}_s$, l'homomorphisme naturel
\[
  \underline{\mathrm{Hom}}(P, \Omega)_s \otimes \mathcal{O}_{S,s}/(\underline{t})
  \longrightarrow
  \mathrm{Hom}_{\mathcal{O}_{S,s}}\bigl(P_s,\ \Omega_s/(\underline{t})\Omega_s\bigr)
\]
soit un isomorphisme\footnote{Le titre, souligné, est lu « Cor.\ Main »,
comme aux pages 5 et 9 du dossier 90 ; ce peut être une abréviation de
« Corollaire ». Il en va de même des corollaires 9, 10 et 11.}.

\emph{Démonstration.} On peut supposer $S = \operatorname{Spec} A$ affine. Une
présentation $L_1 \to L_0 \to P \to 0$ par des libres de type fini fait de
$\mathrm{Hom}(P, \Omega \otimes M)$ le $H^0$ du complexe
$K^{\cdot} \otimes M$, où $K^{\cdot} \colon \check{L}_0 \otimes \Omega \to
\check{L}_1 \otimes \Omega$ est placé en degrés 0 et 1. Si $B^1$ désigne
l'image de $K^0 \to K^1$, la flèche $H^0(K^{\cdot}) \otimes M \to
H^0(K^{\cdot} \otimes M)$ est bijective dès que
$\mathrm{Tor}_1(M, H^1(K^{\cdot})) = \mathrm{Tor}_1(M, B^1) = 0$. Le lemme 6,
appliqué aux Modules cohérents $K^1$ et $H^1(K^{\cdot})$ avec
$M = \mathcal{O}_{S,s}/(\underline{t})$, annule tous leurs $\mathrm{Tor}_i$,
$i > 0$, et la suite $0 \to B^1 \to K^1 \to H^1 \to 0$ annule alors
$\mathrm{Tor}_1(M, B^1)$\footnote{La page demande $\mathrm{Tor}_1(M,
H^1(K^{\cdot})) = \mathrm{Tor}_1(M, Z^1(K^{\cdot})) = 0$, et nomme ces modules
$R$ et $S$ (sans rapport avec le schéma $S$). Pour un complexe concentré en
degrés 0 et 1, $Z^1 = K^1$ ; la condition qui sert est sur $B^1$, et on la
déduit de celle sur $K^1$ par l'annulation de $\mathrm{Tor}_2(M, H^1)$, que le
lemme 6 donne aussi. La lettre écrite comme un rond ouvert dans
$\underline{\mathrm{Hom}}(P, Q)$ est l'$\Omega$ de l'énoncé.}.

\pagerange{10}{12}
\subsection*{Le faisceau fondamental (pages 10 et 12)}

Ce qu'il appelle \emph{faisceau fondamental} sur un schéma de Cohen--Macaulay
$X$ est un Module cohérent $\Omega$ dont chaque fibre $\Omega_x$ est un module
canonique de $\mathcal{O}_{X,x}$ ; le dossier ne le définit pas\footnote{La
définition est celle du « module fondamental » des « Notes antiques 1958 » du
dossier 101 (page 18) : $\Omega$ de type fini dont le dual de Matlis est le
module de cohomologie locale $H^n_{\mathfrak{m}}(A)$, $n = \dim A$ ; une
addition y précise que, pour $A$ de Cohen--Macaulay, module fondamental et
module dualisant coïncident. C'est le module canonique d'aujourd'hui. Le haut
de la page 10 porte une ligne biffée (« Lemme. Soit $X$ un préschéma
régulier ») et, en interligne, « $\Omega$ [\ldots] définissables par un
C.M. » ; le passage n'est lu qu'en partie.}.

\textbf{Lemme 8.} Soient $X$ un schéma de Cohen--Macaulay et $\Omega$ un
faisceau fondamental sur $X$. Pour toute suite $\mathcal{O}_X$-régulière
$\underline{t}$ de sections de $\mathcal{O}_X$, $\Omega/(\underline{t})\Omega$
est un faisceau fondamental sur le schéma de Cohen--Macaulay
$V(\underline{t})$.

C'est le fait classique que la réduction d'un module canonique modulo une
suite régulière est un module canonique du quotient.

\textbf{Corollaire 9.} Pour $x \in X$ et un système de paramètres
$\underline{t}$ de $\mathcal{O}_{X,x}$, $\Omega_x/(\underline{t})\Omega_x$ est
l'enveloppe injective du corps résiduel sur l'anneau artinien
$\mathcal{O}_{X,x}/(\underline{t})$.

En effet $\underline{t}$ est régulière ($\mathcal{O}_{X,x}$ est de
Cohen--Macaulay), et le module canonique d'un anneau local artinien est
l'enveloppe injective de son corps résiduel.

\textbf{Corollaire 10.} Soit $X$ un schéma localement noethérien, localement
immergeable dans un schéma régulier. Il existe un ouvert dense $U$ de $X$ et un
$\mathcal{O}_U$-Module cohérent $\Omega$ tels que :
\begin{enumerate}
\item[a)] les anneaux locaux de $U$ sont de Cohen--Macaulay ;
\item[b)] pour tout $x \in U$ et tout système de paramètres $\underline{t}$ de
$\mathcal{O}_{X,x}$, $\Omega_x/(\underline{t})\Omega_x$ est l'enveloppe
injective du corps résiduel de $\mathcal{O}_{X,x}/(\underline{t})$.
\end{enumerate}
La page ne le démontre pas. Le lieu de Cohen--Macaulay contient les points
maximaux ; si $X$ est fermé de codimension $c$ dans un régulier $Y$, au
voisinage d'un tel point, $\underline{\mathrm{Ext}}^c_{\mathcal{O}_Y}
(\mathcal{O}_X, \mathcal{O}_Y)$ y est un faisceau fondamental, et le
corollaire 9 donne b)\footnote{Cette esquisse est de cette lecture. La page
place $\Omega$ sur $X$ ; seule sa restriction à $U$ sert, et c'est ainsi qu'on
l'énonce. En marge : « il suffit [que $X$] soit excellent, ou seulement qu'il
satisfasse la condition de EGA IV 6.11.8 » ; le § 6.11 d'EGA IV traite de
l'ouverture du lieu de Cohen--Macaulay. Nous ne vérifions pas cette variante,
où l'existence d'un module canonique, et non seulement l'ouverture, est en
jeu. « Localement immergeable dans un régulier » est souligné d'un trait
ondulé.}.

Ces conditions entraînent :
\begin{enumerate}
\item[c)] soient $x \in U$ et $u \colon R \to R'$ un homomorphisme surjectif de
$\mathcal{O}_{X,x}$-modules de type fini ; si, pour tout système de paramètres
$\underline{t}$,
\[
  \mathrm{Hom}\bigl(R',\ \Omega_x/(\underline{t})\Omega_x\bigr) \to
  \mathrm{Hom}\bigl(R,\ \Omega_x/(\underline{t})\Omega_x\bigr)
\]
est bijectif, alors $u$ est un isomorphisme.
\end{enumerate}
En effet, la dualité de Matlis sur l'anneau artinien
$\mathcal{O}_{X,x}/(\underline{t})$ étant fidèle et exacte, $u$ induit un
isomorphisme $R/(\underline{t})R \simeq R'/(\underline{t})R'$, de sorte que
$\ker u \subset (\underline{t})R$ ; prenant les systèmes $(t_1^k, \ldots,
t_n^k)$, $\ker u \subset \bigcap_k \mathfrak{m}_x^k R = 0$ par le théorème de
Krull\footnote{La justification de c) est de cette lecture ; la page dit
« ces conditions impliquent la suivante ». Suit un premier « Cor.\ Main 11 »,
entièrement barré, avec en marge « à expliciter d'une autre façon » : il
formulait la même conclusion à l'aide des foncteurs $\mathrm{Hom}(R, -)$,
$\mathrm{Hom}(R', -)$ et d'une condition de changement de base sur les suites
régulières. La page 14 commence par un second état, barré lui aussi.}.

\pagerange{14}{15}
\subsection*{La stabilisation d'un pro-système (pages 14 et 15)}

\textbf{Corollaire 11.} Soient $U = \operatorname{Spec} A$ et $\Omega$ comme au
corollaire 10, $(P_i)$ un système projectif filtrant strict de $A$-modules de
type fini, et $F = \varinjlim_i \mathrm{Hom}(P_i, -)$ le foncteur qu'il
pro-représente. Supposons :
\begin{enumerate}
\item[(i)] $F(\Omega)$ est un $A$-module de type fini ;
\item[(ii)] pour tout $\mathfrak{p}$ et tout système de paramètres
$\underline{t}$ de $A_{\mathfrak{p}}$, l'application canonique
$F(\Omega) \otimes A_{\mathfrak{p}}/(\underline{t}) \to
F\bigl(\Omega \otimes A_{\mathfrak{p}}/(\underline{t})\bigr)$ est bijective.
\end{enumerate}
Alors le système $(P_i)$ est stationnaire. Précisément, (i) équivaut à ce que
le système inductif des $\mathrm{Hom}(P_i, \Omega)$ soit stationnaire --- ses
transitions sont injectives, le système étant strict, et $A$ est
noethérien ---, et si $i_0$ est tel que $\mathrm{Hom}(P_{i_0}, \Omega) \simeq
\mathrm{Hom}(P_i, \Omega)$ pour $i \geq i_0$, alors $P_i \simeq P_{i_0}$ pour
$i \geq i_0$\footnote{La page dit « un syst.\ projectif de $A$-modules » ;
« strict » et « de type fini » sont ajoutés ici, comme dans le premier état
barré (« syst.\ projectif filtrant strict de Modules coh. ») et comme le donne
le lemme 5 : sans eux, ni l'équivalence ni c) ne s'appliquent.}.

\emph{Démonstration.} Écrivons $E = \Omega_{\mathfrak{p}}/(\underline{t})
\Omega_{\mathfrak{p}}$ et considérons le carré

\begin{tikzcd}[column sep=small, nodes={font=\small}]
\varinjlim_i \mathrm{Hom}(P_i, \Omega) \otimes A_{\mathfrak{p}}/(\underline{t}) \arrow[r, "\sim"] & \varinjlim_i \mathrm{Hom}(P_i, E) \\
\mathrm{Hom}(P_{i_0}, \Omega) \otimes A_{\mathfrak{p}}/(\underline{t}) \arrow[r] \arrow[u, "\sim"] & \mathrm{Hom}(P_{i_0}, E) \arrow[u, "\alpha"']
\end{tikzcd}

où la flèche du haut est (ii) et celle de gauche un isomorphisme par le choix
de $i_0$. Il montre que $\alpha$ est surjectif ; il est injectif, les
transitions du système de droite l'étant. Donc $\mathrm{Hom}(P_{i_0}, E) \to
\mathrm{Hom}(P_i, E)$ est bijectif pour $i \geq i_0$, et c), appliqué à
$(P_i)_{\mathfrak{p}} \to (P_{i_0})_{\mathfrak{p}}$, montre que c'est un
isomorphisme. Ceci valant en tout $\mathfrak{p}$, $P_i \to P_{i_0}$ est
bijectif\footnote{La page conclut que $\alpha$ est un isomorphisme en disant
que le diagramme le « montre » ; la flèche du bas n'a pas besoin d'être
bijective. Elle écrit $\Omega/\underline{f}\,\Omega_{\mathfrak{p}}$ en haut à
droite et $\Omega_{\mathfrak{p}}/\underline{f}\,\Omega_{\mathfrak{p}}$ en bas,
pour le même module.}.

\pagerange{15}{16}
\section*{IV. Le critère (pages 15 et 16)}

« Mettant ensemble prop.~3, lemme~5, Remarque~5 bis, et le corollaire~10, on
trouve » :

\textbf{Théorème 12.} Soient $A$ un anneau noethérien quotient d'un anneau
régulier, et $F \colon \mathrm{Mod}(A) \to (\mathrm{Ens})$ un foncteur. Pour
que $F$ soit représentable par un module de type fini, il faut et il suffit
qu'il satisfasse aux conditions suivantes :
\begin{enumerate}
\item[(i)] $F$ commute aux limites inductives filtrantes ;
\item[(ii)] $F$ est exact à gauche ;
\item[(iii)] a) pour tout $\mathfrak{p}$ et tout
$\widehat{A}_{\mathfrak{p}}$-module de type fini $M$,
$F(M) \xrightarrow{\ \sim\ } \varprojlim F(M/\mathfrak{m}_{\mathfrak{p}}^{n+1}M)$ ;
b) pour tout $g \in A$ et tout $\widehat{A}^{(g)}$-module de type fini $M$,
$F(M) \xrightarrow{\ \sim\ } \varprojlim F(M/g^{n+1}M)$ ;
\item[(iv)] pour tout $A$-module de type fini $M$, $F(M)$ est de type fini ;
\item[(v)] la fonction $\mathfrak{p} \mapsto \dim_{k(\mathfrak{p})}
F(k(\mathfrak{p}))$ est constructible sur $\operatorname{Spec} A$ ;
\item[(vi)] pour tout quotient $B$ de $A$ et tout $B$-module de type fini
$\Omega$, il existe $h \in B$ non nilpotent tel que, pour tout
$\mathfrak{p} \in \operatorname{Spec} B_h$ et tout système de paramètres
$\underline{t}$ de $B_{\mathfrak{p}}$, l'application canonique
\[
  F(\Omega) \otimes B_{\mathfrak{p}}/(\underline{t}) \longrightarrow
  F\bigl(\Omega \otimes B_{\mathfrak{p}}/(\underline{t})\bigr)
\]
soit bijective.
\end{enumerate}
C'est l'exactitude à gauche qui porte le numéro (ii), et la remarque 18 s'y
réfère sous ce numéro\footnote{Sur la page, (ii) est écrit avant (i), le
chiffre de (i) surchargeant un autre, et un trait courbe en marge semble les
intervertir ; on les donne dans l'ordre des numéros. Un mot ajouté au-dessus de
« quotient » en (vi) n'est pas lu.}.

\emph{Nécessité.} Elle est claire pour (i) à (v). Pour (vi), si
$F = \mathrm{Hom}_A(P, -)$, le corollaire 7, appliqué sur
$\operatorname{Spec} B$ à $P \otimes_A B$ et $\Omega$, donne un ouvert dense ;
on y choisit $D(h)$, $h$ non nilpotent, sur lequel $B$ est de
Cohen--Macaulay --- $B$, quotient d'un régulier, a un lieu de Cohen--Macaulay
ouvert, qui contient les points maximaux. Un système de paramètres y est alors
une suite régulière.

\emph{Suffisance.} Les conditions passent de $F$ à $F_B$, $B$ quotient de $A$ ;
par la proposition 3 --- qui utilise (i), (ii), (iii b) ---, il suffit de
trouver $h \in A$ non nilpotent tel que $F_h$ soit représentable. Par le
lemme 5 et la remarque 5 bis --- qui utilisent (i), (ii), (iii a), (iv), (v)
---, $F$ est pro-représentable par un système strict $(P_i)$ de modules de type
fini\footnote{La condition (iv) du lemme 5, sur les modules de longueur finie
sur $A_{\mathfrak{p}}$, résulte de (iv) ici : un tel module est
$M_0 \otimes A_{\mathfrak{p}}$ pour un $A$-module de type fini $M_0$, et
$F(M_0 \otimes A_{\mathfrak{p}}) = F(M_0)_{\mathfrak{p}}$. La remarque 5 bis
demande que le lieu normal de chaque quotient intègre de $A$ contienne un
ouvert non vide ; la page ne le vérifie pas pour un quotient d'un anneau
régulier, et nous ne savons pas si c'est toujours le cas. C'est vrai si $A$
est excellent, ce qui couvre l'application du théorème 14.}. Le corollaire 10
fournit un ouvert dense de $\operatorname{Spec} A$, qu'on peut prendre de la
forme $D(g)$ et dense, et un faisceau fondamental sur lui, qu'on prolonge en un
$A$-module de type fini $\Omega$. La condition (vi), appliquée à $B = A$ et à
$\Omega$, donne $h'$ ; $D(gh') = D(g) \cap D(h')$ est non vide, $D(g)$ étant
dense. Sur $\operatorname{Spec} A_{gh'}$, le corollaire 11 s'applique au
système $(P_{i, gh'})$ --- sa condition (i) par (iv), sa condition (ii) par
(vi) ---, qui est donc stationnaire : $F_{gh'}$ est représentable, cqfd\footnote{En marge,
en face des trois phrases de la démonstration, les groupes de conditions
« (i) (ii) (iii b) (iv) », « (i) (ii) (iii a) (iv) (v) », « (vi) ». La page
biffe une récurrence noethérienne, inutile puisque la proposition 3 la
contient. Le choix de $D(g)$ dense et du produit $gh'$ est de cette lecture ;
la page dit « on peut maintenant [trouver] $f$, non nilpotent, tel que le
syst.\ $P_{i,f}$ soit stationnaire ».}.

\textbf{Remarque 13.} Supposons que tout quotient intègre $B$ de $A$ ait un
ouvert non vide $D(h)$ avec $B_h$ régulier. Selon la page, cette condition
rend inutile l'hypothèse que $A$ soit quotient d'un régulier, et entraîne
que, dans le théorème 12, (vi) implique (v) ; la page s'interrompt
là\footnote{La phrase se termine, biffée, par « et dans (vi) », au bas de la
page 16 ; la page 17 commence le théorème 14.}.

La seconde affirmation se justifie ainsi. Soit $\mathfrak{p}$ premier,
$B = A/\mathfrak{p}$. La condition (vi) pour $B$ et $\Omega = B$, et
l'hypothèse, donnent un ouvert non vide de $\operatorname{Spec} B$ où $B$ est
régulier et où $F(B) \otimes B_{\mathfrak{q}}/(\underline{t}) \simeq
F(B_{\mathfrak{q}}/(\underline{t}))$ ; en prenant pour $\underline{t}$ un
système régulier de paramètres, $B_{\mathfrak{q}}/(\underline{t}) =
k(\mathfrak{q})$, et $\dim F(k(\mathfrak{q}))$ est le rang de la fibre du
module de type fini $F(B)$, constant sur un ouvert non vide. C'est le critère
de constructibilité rappelé à la page 7. La première affirmation est claire
quand $A$ est intègre --- sur l'ouvert régulier, $\Omega = \mathcal{O}$ est un
faisceau fondamental et remplace le corollaire 10 ---, mais la page ne dit pas
comment traiter un quotient non réduit, et nous la laissons
ouverte\footnote{Les deux justifications sont de cette lecture.}.

\pagerange{17}{20}
\section*{V. Le module $P$ d'un morphisme propre (pages 17 à 20)}

\pagerange{17}{18}
\subsection*{L'énoncé (pages 17 et 18)}

\textbf{Théorème 14.} Soient $f \colon X \to S$ un morphisme propre et de
présentation finie, $\mathcal{F}$ et $\mathcal{G}$ des $\mathcal{O}_X$-Modules
de présentation finie, $\mathcal{G}$ plat sur $S$. Il existe un
$\mathcal{O}_S$-Module de présentation finie $P$ et un isomorphisme,
fonctoriel en le $\mathcal{O}_S$-Module quasi-cohérent $M$,
\[
  \underline{\mathrm{Hom}}_{\mathcal{O}_S}(P, M) \simeq
  f_{*}\,\underline{\mathrm{Hom}}_{\mathcal{O}_X}(\mathcal{F},\ \mathcal{G}
  \otimes_{\mathcal{O}_S} M) .
\]
En prenant les sections globales, $\mathrm{Hom}(P, M) \simeq
\mathrm{Hom}(\mathcal{F}, \mathcal{G} \otimes_{\mathcal{O}_S} M)$, ce qui
détermine $P$ à isomorphisme unique près. La formation de $P$ commute à tout
changement de base $S' \to S$\footnote{C'est immédiat pour $S'$ affine : les
deux membres, pour un $\mathcal{O}_{S'}$-Module $M'$, se calculent sur $X$ en
regardant $M'$ comme $\mathcal{O}_S$-Module. Le théorème se trouve, pour une
base localement noethérienne et des Modules cohérents, au § 7.7 d'EGA III
(seconde partie, 1963), par une autre démonstration, à partir d'un complexe
qui calcule la cohomologie. En marge gauche, une note oblique, lue en partie :
« \emph{NB}. Il y a le même énoncé de représentabilité, par un $W(\ldots)$
[\ldots], pour $F$ de prés.\ finie, $G$ quasi-cohérent ; [\ldots] représentable
par un $P$ [\ldots] de prés.\ finie [\ldots] ». Nous ne la complétons pas.}.

On se ramène à $S = \operatorname{Spec} A$ affine, puis, comme d'habitude
--- tout étant de présentation finie, et $\mathcal{G}$ plat ---, à $A$ de type
fini sur $\mathbb{Z}$, donc quotient d'un anneau régulier, et l'on applique le
théorème 12 au foncteur $M \mapsto \mathrm{Hom}(\mathcal{F}, \mathcal{G}
\otimes_A M)$. Les conditions (i) à (iv) sont « standard » : (i) parce que $X$
est quasi-compact et $\mathcal{F}$ de présentation finie, (ii) par la
platitude de $\mathcal{G}$, (iii) par le théorème de comparaison (fonctions
formelles) d'EGA III, (iv) par le théorème de finitude pour un morphisme
propre. $A$ étant excellent, la remarque 13, en sa partie justifiée ci-dessus,
réduit (v) à (vi)\footnote{Un passage barré commençait la vérification directe
de (v) : « $A$ intègre, il existe un ouvert non vide $U$ \ldots\ tel que sur
$U$, la fonction $s \mapsto \dim \mathrm{Hom}(F_s, G_s)$ soit constante » ;
« En vertu de la remarque 13, il reste à vérifier la condition (vi) » est écrit
au-dessus.}. Il reste (vi).

Sauf l'exactitude à gauche, les conditions du théorème 12 se vérifient en fait
pour tous les foncteurs
\[
  M \longmapsto H^{q}\bigl(X,\ \underline{\mathrm{Ext}}^{i}(\mathcal{F},\
  \mathcal{G} \otimes_A M)\bigr)
  \qquad\text{ou}\qquad
  M \longmapsto \mathrm{Ext}^{q}(X;\ \mathcal{F},\ \mathcal{G} \otimes_A M),
\]
sans supposer $\mathcal{G}$ plat --- sur une base quelconque, $M \mapsto
R^q f_{*}\,\underline{\mathrm{Ext}}^{i}(\mathcal{F}, \mathcal{G} \otimes M)$
et $M \mapsto R^q f_{*}\,R\underline{\mathrm{Hom}}(\mathcal{F}, \mathcal{G}
\otimes M)$. La vérification se fait « en deux coups » : changement de base
générique pour les $\underline{\mathrm{Ext}}$ (lemme 15), puis pour les $R^j
f_{*}$ (lemme 17). C'est d'ici que part la remarque 18.

\pagerange{18}{20}
\subsection*{Changement de base générique (pages 18 à 20)}

\textbf{Lemme 15.} Soient $A$ noethérien, $f \colon X \to S =
\operatorname{Spec} A$ de type fini, $\mathcal{F}$, $\mathcal{G}$ cohérents sur
$X$, $\Omega$ cohérent sur $S$, $N$ un entier. Il existe un ouvert dense $U$ de
$S$ tel que, pour $\mathfrak{p} \in U$ et toute suite
$A_{\mathfrak{p}}$-régulière $\underline{t}$,
\[
  \underline{\mathrm{Ext}}^{i}(\mathcal{F},\ \mathcal{G} \otimes_A \Omega)
  \otimes A_{\mathfrak{p}}/(\underline{t})
  \xrightarrow{\ \sim\ }
  \underline{\mathrm{Ext}}^{i}\bigl(\mathcal{F},\ \mathcal{G} \otimes_A \Omega
  \otimes A_{\mathfrak{p}}/(\underline{t})\bigr)
  \qquad (i \leq N) .
\]

\emph{Démonstration.} On peut supposer $X = \operatorname{Spec} B$ affine
($X$ est quasi-compact). Une résolution $L_{\cdot}$ de $\mathcal{F}$ par des
$B$-modules libres de type fini donne, avec $\mathcal{G}' = \mathcal{G}
\otimes_A \Omega$ et $K^{\cdot} = \check{L}_{\cdot} \otimes_B \mathcal{G}'$,
\[
  \underline{\mathrm{Ext}}^{i}(\mathcal{F}, \mathcal{G}' \otimes_A M) =
  H^{i}(K^{\cdot} \otimes_A M)
\]
pour tout $A$-module $M$, et il s'agit que $H^{i}(K^{\cdot}) \otimes_A M \to
H^{i}(K^{\cdot} \otimes_A M)$ soit bijectif pour $i \leq N$. Quitte à tronquer
$K^{\cdot}$ au-delà du degré $N+1$, il est borné, et il suffit que
$\mathrm{Tor}_i^A(K^{j}, M) = \mathrm{Tor}_i^A(H^{j}(K^{\cdot}), M) = 0$ pour
$i > 0$ et tout $j$. Les $K^j$ et $H^j(K^{\cdot})$ sont des $B$-modules de type
fini --- des Modules cohérents sur $X$, non sur $S$ ---, et l'on est ramené
au

\textbf{Lemme 16.} Soit $R$ un Module cohérent sur $X$. Il existe un ouvert
dense $U$ de $S$ tel que, pour $\mathfrak{p} \in U$ et toute suite
$A_{\mathfrak{p}}$-régulière $\underline{t}$,
$\mathrm{Tor}_i^{A}(R,\ A_{\mathfrak{p}}/(\underline{t})) = 0$ pour $i > 0$.

\emph{Démonstration.} D'après celle du lemme 6, il suffit que, $J$ étant le
nilradical de $A$, $\mathrm{gr}_J(A)$ et $\mathrm{gr}_J(R)$ soient plats sur
$A/J$ aux points de $U$ ; le théorème de platitude générique permet de choisir
$U$ ainsi\footnote{La page écrit $\mathrm{gr}_J(\mathcal{O}_X)$ là où le
raisonnement demande $\mathrm{gr}_J(R)$, et garde $F$ dans la formule après
avoir ajouté en interligne « Soit $R$ cohérent sur $X$ ». Dans le lemme 6 la
liberté des gradués servait deux fois : à voir que $\bar{\underline{t}}$ est
régulière dans $A/J$ (ce qui ne fait intervenir que $A$), et à transporter la
régularité au module, pour quoi la platitude suffit. Un module plat n'étant pas
de type fini sur $A$, on n'utilise pas Nakayama ; il n'en est pas besoin.}.

La fin de la page 19 est raturée et inachevée : le lemme 15 ramène le cas de
$R^q f_{*}\,\underline{\mathrm{Ext}}^{i}(\mathcal{F}, \mathcal{G} \otimes
\Omega)$ au

\textbf{Lemme 17.} Soient $S$ noethérien, $f \colon X \to S$ propre et $E$
cohérent sur $X$. Il existe un ouvert dense $U$ de $S$ tel que, pour tout
$s \in U$ et toute suite $\mathcal{O}_{S,s}$-régulière $\underline{t}$,
l'homomorphisme naturel
\[
  R^{j} f_{*}(E)_s \otimes \mathcal{O}_{S,s}/(\underline{t}) \longrightarrow
  R^{j} f^{(s)}_{*}\bigl(E^{(s)} \otimes \mathcal{O}_{S,s}/(\underline{t})\bigr)
\]
soit un isomorphisme pour tout $j$, où $f^{(s)} \colon X^{(s)} = X \times_S
\operatorname{Spec} \mathcal{O}_{S,s} \to \operatorname{Spec}
\mathcal{O}_{S,s}$ et $E^{(s)}$ est induit par $E$.

\emph{Démonstration.} Posons $M = \mathcal{O}_{S,s}/(\underline{t})$. Par le
lemme 16, on peut supposer $E \otimes M = E \otimes^{\mathbb{L}} M$. On
représente $Rf_{*}E$, localement sur $S$, par un complexe borné $K^{\cdot}$ à
composantes cohérentes, et le lemme 16, appliqué avec $X = S$, permet de
supposer $K^{\cdot} \otimes M = K^{\cdot} \otimes^{\mathbb{L}} M$. La formule
de projection
\[
  Rf_{*}(E \otimes^{\mathbb{L}} M) \xrightarrow{\ \sim\ }
  Rf_{*}(E) \otimes^{\mathbb{L}} M
\]
s'identifie alors à $H^{j}(K^{\cdot}) \otimes M \to H^{j}(K^{\cdot} \otimes
M)$ en chaque degré. Restreignant encore $U$ par le lemme 16, on peut supposer
$\mathrm{Tor}_i(H^{j}(K^{\cdot}), M) = 0$ pour $i > 0$ et tout $j$, et cette
flèche est alors un isomorphisme, cqfd\footnote{La page appelle la flèche
$\alpha$ « l'isom.\ de régularité », lecture incertaine ; c'est la formule de
projection. La lecture de la seconde accolade
($K^{\cdot} \otimes^{\mathbb{L}} M = K^{\cdot} \otimes M$) est incertaine. La
page écrit $M = A_{\mathfrak{p}}/\underline{f}A_{\mathfrak{p}}$, en interligne.}.

Mises ensemble, la platitude de $\mathcal{G}$ et les lemmes 15 et 17
(appliqués à $E = \underline{\mathrm{Hom}}(\mathcal{F}, \mathcal{G}_B \otimes
\Omega)$ sur $X_B$, avec $N = 0$) donnent (vi) pour le foncteur du théorème
14 ; le théorème 12 conclut. La page ne rédige pas ce dernier assemblage.

\pagerange{21}{23}
\section*{VI. Foncteurs constructibles, restriction de Weil (pages 21 à 23)}

\pagerange{21}{21}
\subsection*{Foncteurs constructibles (page 21)}

\textbf{Remarque 18.} Il y a lieu d'appeler \emph{foncteur constructible} un
foncteur qui satisfait aux conditions du théorème 12, sauf l'exactitude à
gauche (ii). Il faudrait définir de même la notion de foncteur constructible
$\mathrm{Qcoh}(Y) \to \mathrm{Qcoh}(X)$ (pour $X$ au-dessus de $Y$), et il y
aurait une suite : un composé de foncteurs constructibles l'est encore, et les
foncteurs $\underline{\mathrm{Ext}}^{i}$, $\underline{\mathrm{Tor}}_{j}$,
$R^{q}f_{*}$ sont constructibles\footnote{Au crayon, dans l'angle, « (additif
?) », relié à « foncteur » : sans exactitude à gauche, un foncteur à valeurs
dans les ensembles n'a pas de raison d'être additif, et la notion demande
sans doute qu'on le suppose. La lecture de l'exposant $q$ de $R^q f_{*}$ est
incertaine ; la phrase s'achève sur des points de suspension. Le programme
s'apparente --- sans s'y identifier, les conditions étant d'une autre nature
--- aux foncteurs \emph{cohérents} d'Auslander (« Coherent functors », 1966),
que Hartshorne a repris pour les modules sur un anneau noethérien (« Coherent
functors », \emph{Advances in Math.}\ 140, 1998) en montrant que les foncteurs
de cohomologie d'un morphisme propre le sont ; ils servent aujourd'hui à
démontrer des énoncés comme le théorème 14 sur des bases très générales. Le
dossier ne pouvait pas connaître ces travaux, sauf peut-être le premier.}.

\pagerange{21}{23}
\subsection*{La restriction de Weil (pages 21 à 23)}

\textbf{Théorème 19.} Soit $f \colon X \to S$ un morphisme propre, plat et de
présentation finie. Pour tout schéma $Z$ \emph{affine} sur $X$, le foncteur
\[
  f_{*}(Z) = \prod_{X/S} Z/X, \qquad
  T \longmapsto \mathrm{Hom}_X(X \times_S T,\ Z),
\]
est représentable par un schéma affine sur $S$ ; si $Z$ est de présentation
finie (resp.\ de type fini) sur $X$, $f_{*}(Z)$ l'est sur $S$\footnote{C'est
la \emph{restriction de Weil} de $Z$ le long de $f$ ; la description par
$T \mapsto \mathrm{Hom}_X(X_T, Z)$ est de cette lecture, la page écrit
seulement le produit, $X/S$ en indice. Le dossier 100, aux pages 47 à 50, s'en
sert le long d'une extension finie et plate de traits, cas où la
représentabilité est classique ; ici la base est propre et plate, et le
résultat est le même pour $Z$ affine. « Affine » est souligné trois fois.}.

\emph{Démonstration.} On peut supposer $S$ affine, donc $X$ quasi-compact. Alors
$Z = \varprojlim Z_i$, limite filtrante de $X$-schémas affines de présentation
finie, et $f_{*}(Z) = \varprojlim f_{*}(Z_i)$ : pour la première assertion, il
suffit de traiter les $Z_i$, une limite de $S$-schémas affines étant affine.
Pour $Z$ de type fini, on peut prendre les transitions $Z_j \to Z_i$ et les
$Z \to Z_i$ des immersions fermées ; s'il est prouvé que les $f_{*}(Z_j) \to
f_{*}(Z_i)$ le sont aussi, $f_{*}(Z) \to f_{*}(Z_i)$ est une immersion fermée
et $f_{*}(Z)$ est de type fini. On est ramené à prouver en même temps le

\textbf{Corollaire 20.} Si $Z \to Z'$ est une immersion fermée de $X$-schémas
affines, $f_{*}(Z) \to f_{*}(Z')$ est une immersion fermée.

Par les arguments de limite standard, on se ramène à $Z$, $Z'$ de présentation
finie, puis à $S$ de type fini sur $\mathbb{Z}$, donc noethérien.

\emph{Premier cas : $Z = \mathbb{V}(\mathcal{F})$}, fibré vectoriel
$\operatorname{Spec} \mathrm{Sym}(\mathcal{F})$ d'un $\mathcal{F}$ cohérent sur
$X$. Le théorème 14, avec $\mathcal{G} = \mathcal{O}_X$ --- plat sur $S$
puisque $X$ l'est ---, fournit $P$ cohérent sur $S$ avec
$\mathrm{Hom}(\mathcal{F}, f^{*}M) \simeq \mathrm{Hom}(P, M)$, et l'on voit
aussitôt que $f_{*}(\mathbb{V}(\mathcal{F})) \simeq \mathbb{V}(P)$ : une
section de $\mathbb{V}(\mathcal{F})$ sur $X_T$ est un homomorphisme
$\mathcal{F}_T \to \mathcal{O}_{X_T}$\footnote{La page invoque « le th.~12 » ;
c'est le théorème 14, qui en est la conséquence, dans le cas
$\mathcal{G} = \mathcal{O}_X$.}.

\emph{Second cas : le cas général.} On a une immersion fermée
$Z \hookrightarrow Z' = \mathbb{V}(\mathcal{F})$, et il suffit de prouver que
$f_{*}(Z) \to f_{*}(Z')$ est représentable par des immersions fermées ---
ce qui prouve du même coup le corollaire 20. Un $T$-point de $f_{*}(Z')$ est
une section $s$ de $Z'_T$ sur $X_T$ ; le sous-foncteur de $T$ qui « fait passer
$s$ par $Z$ » est $f_{*}(Y) = \prod_{X_T/T} Y/X_T$, où $Y = s^{-1}(Z_T)$ est un
sous-schéma fermé de $X_T$. Il est représentable par un sous-schéma fermé de
$T$, par « l'exposé Murre, th.~3 », cqfd\footnote{La lecture « Murre » est
incertaine. Il s'agit vraisemblablement de l'exposé de J.~P. Murre au
Séminaire Bourbaki, « Representation of unramified functors. Applications
(according to unpublished results of A. Grothendieck) » (exposé 294, 1965) ;
nous n'avons pas vérifié que son théorème 3 soit l'énoncé invoqué. On peut
aussi le tirer du théorème 14 : si $\mathcal{I}$ est l'idéal de $Y$, la
condition $Y_{T'} = X_{T'}$ est l'annulation de $\mathcal{I} \to
\mathcal{O}_{X} \otimes \mathcal{O}_{T'}$, c'est-à-dire d'une flèche
$P \to \mathcal{O}_{T'}$, et définit le sous-schéma fermé de $T$ d'idéal
l'image de $P$. La page dessine en marge le schéma $Z \hookrightarrow Z'$
au-dessus de $X$, avec la section $s$ de $X$ vers $Z'$, sans la nommer.}.

\textbf{Remarque 21.} Soit $f \colon X \to S$ propre, plat, de présentation
finie. Pour tout $\mathcal{O}_X$-Module quasi-cohérent $\mathcal{F}$, il existe
un $\mathcal{O}_S$-Module quasi-cohérent $P$ tel que
$f_{*}(\mathbb{V}(\mathcal{F})) \simeq \mathbb{V}(P)$ : on écrit $\mathcal{F}$
comme limite inductive filtrante de Modules de présentation finie
$\mathcal{F}_i$, les $P_i$ correspondants forment un système inductif, et
$P = \varinjlim P_i$ convient, puisque $\mathbb{V}(\varinjlim P_i) =
\varprojlim \mathbb{V}(P_i)$\footnote{La page dit seulement « on procède par
passage à la limite à partir du cas $F$ de prés.\ finie » ; elle écrit
$V(F)$ pour le fibré vectoriel, que l'on note $\mathbb{V}$.}.

\end{document}
